\section{模型检验}

本章在第四部分“模型的建立与求解”的基础上，对任务一至任务四的模型输出进行一致性检验。第四部分已经给出主要参数、结果表和图形，本章不再重复列示相同表格，而是从物理一致性、统计有效性、跨域可迁移性和工程可执行性四个角度说明模型结果是否可信。

\subsection{任务一模型检验：飞轮退化拟合与 RUL 合理性}

任务一的检验重点有两个：一是物理指数模型能否解释附件 1 的全寿命电流退化，二是附件 2 当前阶段与 RUL 预测是否相互一致。由表~\ref{tab:q1_solution_params} 可知，指数模型的 RMSE 为 $0.0077\ \mathrm A$，拟合优度 $R^2=0.9997$，说明“润滑退化导致摩擦力矩上升，进而引起电流升高”的机理假设能够较好解释全寿命轨迹。由于指数模型在健康期、退化期和后期加速段均未出现明显系统偏差，其作为飞轮目标域基准模型是合理的。

附件 2 的在线评估结果见表~\ref{tab:q1_online_solution}。确定性指数模型给出 RUL 点估计 $1548.5$ 天，Wiener 概率模型给出均值 $1000.5$ 天，二者均表明当前飞轮尚未达到电流硬失效阈值。与此同时，当前电流为 $1.0088\ \mathrm A$，标准化退化进度约为 $0.646$，已超过衰退阈值 $\theta_F=0.50$，因此阶段判定为 $F$。这说明“尚未立即失效”和“已进入衰退阶段”并不矛盾，前者描述寿命距离，后者描述健康裕度消耗程度。

\subsection{任务二模型检验：轴承特征与健康指数有效性}

任务二需要检验所选振动特征是否真正反映轴承退化。表~\ref{tab:q2_feature_solution} 中，RMS、PeakToPeak 和 MaxFreqAmp 的综合评分排名前三，分别对应振动能量抬升、冲击幅值增强和频谱响应变化。三类特征的物理含义与滚动轴承损伤由局部冲击逐渐演化为宽带振动增强的过程一致，因此可以作为健康指数 HI 的输入。

由表~\ref{tab:q2_hi_solution} 可知，不同源轴承的寿命长度和 HI 拟合效果存在明显差异。Bearing3\_5 的 Wiener 拟合效果最好，$R^2=0.8387$、RMSE 为 $0.0972$，说明其 HI 轨迹更接近稳定的单调退化过程；Bearing3\_1 虽寿命较长，但线性 Wiener 拟合效果较差，提示其退化过程具有更强阶段性。该结果支持本文后续不直接迁移绝对寿命，而采用归一化曲率和域差异筛选的处理方式。

\subsection{任务三模型检验：迁移可行性与不确定性收缩}

任务三的关键是避免无条件迁移。表~\ref{tab:q3_domain_solution} 显示，Bearing3\_5 的 MMD 和 CORAL 均最小，分别为 $0.2589$ 和 $1.4520\times10^{-4}$，满足 $MMD<0.3$ 且 $CORAL<0.001$ 的可迁移判据；Bearing3\_1 的两项域差异指标均偏大，不宜作为主要迁移源。因此，迁移源选择由归一化退化形态的分布差异决定，而非主观指定。

迁移后的 RUL 结果见表~\ref{tab:q3_tl_solution}。迁移模型没有改变附件 2 的 RUL 中心估计，均值仍为 $1000.5$ 天，但将 $95\%$ 置信区间收缩至 $[653.6,1426.3]$ 天。与任务一基准 Wiener 模型相比，迁移学习主要起到约束外推不确定性的作用，而不是强行覆盖目标域观测趋势。因此，任务三符合“目标域观测决定中心趋势，源域形态补充全寿命先验”的设计预期，未出现负迁移迹象。

\subsection{任务四模型检验：预警等级与健康管理一致性}

任务四检验预警机制是否符合工程逻辑。附件 2 第 $1800$ 天的退化进度约为 $0.646$，阶段标签为 $F$。按照表~\ref{tab:warning_rules} 的规则，$S=F$ 至少应触发 $L_2$ 以上等级，而 $L_3$ 的触发条件中包含“$\hat S=F$”，因此当前预警等级判定为 $L_3$。

从概率角度看，迁移模型下短期失效概率较低，这与 RUL 均值仍约为 $1000.5$ 天是一致的。其含义是：飞轮已经处于健康裕度消耗较大的衰退阶段，但尚未逼近立即失效窗口。本文采用“阶段约束优先、概率信息辅助”的保守策略，能够避免因短期概率较低而低估衰退阶段风险。表~\ref{tab:q4_report_solution} 和表~\ref{tab:q4_action_solution} 给出的健康管理建议与该判定一致，具备工程可解释性和可执行性。

\subsection{综合一致性检验}

综合来看，四项任务之间没有出现方向性矛盾。任务一证明飞轮电流退化模型具有较好拟合效果；任务二证明轴承 HI 能够提取源域全寿命退化形态；任务三证明迁移学习能够在不改变目标域中心趋势的前提下收缩 RUL 不确定性；任务四则将阶段标签、RUL 分布和短期失效概率转化为保守预警等级。

进一步结合第六章的灵敏度分析可知，阶段阈值、失效电流、Wiener 扩散强度、迁移带宽和预警阈值的扰动会影响局部数值和区间宽度，但并未改变附件 2 当前处于 $F$ 阶段并触发 $L_3$ 预警的核心结论。因此，本文模型在退化机理、统计拟合、迁移约束和工程预警四个层面具有较好的一致性与稳健性。
